\documentclass[11pt]{article}
\usepackage[a4paper,margin=2.2cm]{geometry}
\usepackage{amsmath,amssymb}
\usepackage{booktabs,longtable,array,tabularx}
\usepackage[T1]{fontenc}
\usepackage{lmodern}
\usepackage{microtype}
\usepackage{graphicx}
\usepackage[strings]{underscore}
\usepackage{float}
\usepackage{xcolor}
\usepackage[hidelinks]{hyperref}
\sloppy
\setlength{\parskip}{0.45em}
\setlength{\parindent}{0pt}
\setcounter{tocdepth}{1}
\renewcommand{\arraystretch}{1.08}
\newcolumntype{L}{>{\raggedright\arraybackslash}X}
\newcommand{\bros}{\beta_r}
\newcommand{\gap}{\Delta E_h(0)^{*}}

\usepackage{xr}
\externaldocument[F-]{figure_methods}
\newcommand{\compname}{figure methods}
\title{Methods}
\author{}
\date{}
\begin{document}
\maketitle
\tableofcontents

\section*{Conventions}
\addcontentsline{toc}{section}{Conventions}


\emph{Companion document.} How each figure of the dossier was generated -- its
population, selection, solver settings, plotted quantities and the numbers in its
caption -- is described in \texttt{figure\_methods.pdf}, one section per figure,
titled by the index the figure carries in the manuscript, the Extended Data or the
Supplementary Information. Section, table and equation numbers given below as ``of
the figure methods'' are numbers in that document.

\emph{Conventions.} \emph{Tregs OFF} and \emph{Tregs ON} denote the two Treg
conditions of one parameter combination (Section~\ref{sec:model}); in Fig.~2
alone, ON and OFF refer to efferocytosis intact and knocked out. An \emph{initial
condition} is a state at $t = 0$; a \emph{final state} is the state at the end of an
integration (not a time average, unless a tail mean is stated). All
model quantities are dimensionless state units (Section~\ref{sec:units}); time is in model
time units. Large integers carry thousands separators; percentages are of the
population named in the same sentence. Classes of parameter combinations are
written \emph{bistable}, \emph{ON\_mono}, \emph{mono\_A} and
\emph{unresolved\_on}; branches are \emph{sterile}, \emph{pathogenic help} and
\emph{pathogenic harm} (Section~\ref{sec:classes}). Scripts are named relative to
the repository root.
\section{Model}
\label{sec:model}

%% ===========================================================================
%% core.tex -- the model, the parameter sweep, the gate, and Figure 3 with Supplementary
%% Figures S01, S06.
%% Body text only (no preamble). Requires amsmath, booktabs.
%% Every number below was recomputed on 2026-10-01 from the data files by
%%   FP_DIR=<folder holding fp_big_batch_*.rds> Rscript pipeline/census_core.R
%% ("census" below; the tag in square brackets is the section of its output), or by
%% the figure scripts named in place. Code references are to the repository.
%% ===========================================================================

%% ---- from docs/methods_parts/core.tex, lines 15-80
\subsection{State variables}

The model is a system of 13 ordinary differential equations for two compartments,
the gut lumen and the lamina propria (LP), separated by the epithelium. The states
are: luminal pathogens $P_L$; LP pathogens $P_{LP}$ and LP commensals $C_{LP}$;
healthy, apoptotic and necrotic epithelial cells $E_h$, $E_a$, $E_n$; resting and
activated (M1) macrophages $M_0$, $M_1$; reactive oxygen species $R$; resting and
active regulatory T cells $T_0$, $T_r$; a danger signal $D$; and a suppressive
signal $S$. Luminal commensals are a fixed reservoir, $C_L = 1000$, not a state.
Macrophages and Tregs are conserved pools: $dM_0/dt = -dM_1/dt$ and
$dT_0/dt = -dT_r/dt$, so that $M_0 + M_1 = K_M$ and $T_0 + T_r = K_T$ for all $t$.

\subsection{Equations}

With $\theta \in \{0,1\}$ switching Treg induction off or on
(\texttt{rhs\_app} in \texttt{pipeline/01\_sweep.R} and \texttt{lib/skeleton\_defs.R};
the two are the same function, apart from the knockout factor of
Section~\ref{sec:model-mods}, which equals 1 everywhere but in Fig.~2),
\begin{align}
\frac{dP_L}{dt} &= r_P\,P_L\Big(1-\frac{P_L}{K_{L,P}}\Big) - \alpha_{\mathrm{AMP}}\,D\,P_L,\\
\frac{dP_{LP}}{dt} &= \gamma_{B,P}\,P_L\Big(1-\frac{E_h}{K_E}\Big)\Big(1-\frac{P_{LP}}{K_{LP,P}}\Big)
  + r_P\,P_{LP}\Big(1-\frac{P_{LP}}{K_{LP,P}}\Big)\nonumber\\
  &\qquad - P_{LP}\big(\alpha_0 M_0 + \alpha_{1,e} M_1 + \alpha_{1,r} R\big),\\
\frac{dC_{LP}}{dt} &= \gamma_{B,C}\,C_L\Big(1-\frac{E_h}{K_E}\Big)\Big(1-\frac{C_{LP}}{K_{LP,C}}\Big)
  - C_{LP}\big(\alpha_0 M_0 + \alpha_{1,e} M_1 + \alpha_{1,r} R\big),\\
\frac{dE_h}{dt} &= \kappa_E\,(K_E - E_h) - E_h\Big(\tau_r R + \tau_{P,B}\,\frac{P_{LP}}{K_{LP,P}}\Big),\\
\frac{dE_a}{dt} &= \tau_r\,R\,E_h - \mu_{E_a} E_a - g(L)\,(M_0+M_1)\,E_a,\\
\frac{dE_n}{dt} &= \tau_{P,B}\,\frac{P_{LP}}{K_{LP,P}}\,E_h + \mu_{E_a} E_a - g(L)\,(M_0+M_1)\,E_n,\\
\frac{dM_0}{dt} &= -\tau_{01}\,D\,M_0 + \tau_{10}\,S\,M_1 + \tau_{1,d}\,M_1 + \tau_{\mathrm{eff}}\,g(L)\,E_a M_1,\\
\frac{dM_1}{dt} &= \tau_{01}\,D\,M_0 - \tau_{10}\,S\,M_1 - \tau_{1,d}\,M_1 - \tau_{\mathrm{eff}}\,g(L)\,E_a M_1,\\
\frac{dR}{dt} &= \beta_r\,M_1 - \mu_R\,R,\\
\frac{dT_0}{dt} &= -\theta\,\tau_{t,a}\,M_1\,C_{\mathrm{rat}}\,T_0 + \tau_{t,d}\,T_r,\\
\frac{dT_r}{dt} &= \theta\,\tau_{t,a}\,M_1\,C_{\mathrm{rat}}\,T_0 - \tau_{t,d}\,T_r,\\
\frac{dD}{dt} &= \sigma_D\big(w\,f_{E_n} + (1-w)\,f_{PC}\big) - \tau_D\,D,\\
\frac{dS}{dt} &= \sigma_S\,T_r - \tau_S\,S,
\label{eq:model-last}
\end{align}
with
\begin{gather}
L = E_a + E_n,\qquad
g(L) = a_{\mathrm{floor}} + (a_0 - a_{\mathrm{floor}})\,\frac{K_g^{\,n_g}}{K_g^{\,n_g} + L^{n_g}},
\label{eq:g}\\
f_P = \frac{P_{LP}}{K_{LP,P}},\qquad f_C = \frac{C_{LP}}{K_{LP,C}},\qquad
C_{\mathrm{rat}} = \begin{cases} f_C/(f_P+f_C), & f_P + f_C > 0,\\ 0, & \text{otherwise},\end{cases}\nonumber\\
f_{PC} = \frac{P_{LP} + C_{LP}}{K_{LP,P} + K_{LP,C}},\qquad
f_{E_n} = \frac{E_n}{K_{E_n} + E_n},\nonumber\\
K_{E_n} = \frac{K_E\,\tau_{P,B} + \mu_{E_a} K_{E_a}}{a_0 K_M},\qquad
K_{E_a} = \frac{K_E\,\tau_r\,K_R}{\mu_{E_a} + a_0 K_M}.
\label{eq:KEaKEn}
\end{gather}
A single pathogen growth rate $r_P$ applies in the lumen and in the LP (code
parameters \texttt{r\_L\_P} and \texttt{r\_LP\_P}, set equal in every call).
Efferocytic clearance is the decreasing Hill function $g(L)$ of the dead-cell load:
close to $a_0$ for $L \ll K_g$, equal to $(a_0 + a_{\mathrm{floor}})/2$ at $L = K_g$, and
approaching $a_{\mathrm{floor}}$ for $L \gg K_g$. The same factor multiplies the
efferocytosis-driven reversion of M1 macrophages, $\tau_{\mathrm{eff}}\,g(L)\,E_a M_1$.
$K_{E_a}$ and $K_{E_n}$ are normalizing constants built from the intact $a_0$.
The right-hand side is evaluated at the componentwise positive part of the state,
$\max(y, 0)$; for trajectories in the non-negative orthant this changes nothing, and
it prevents integration round-off below the invariant manifold $P_{LP} = 0$ from
engaging the logistic term, which is negative and quadratic for $P_{LP} < 0$.

\subsection{Treg conditions}
\label{sec:treg-conditions}

\emph{Tregs OFF} is $\theta = 0$: $T_r$ and $S$ start at zero and stay there, so
every Treg term vanishes. \emph{Tregs ON} is $\theta = 1$. Nothing
else differs between the two conditions of a parameter combination. Tregs are
the \emph{baseline} Treg, fully stable ($\phi = 0$) and fully specific ($\eta = 1$),
in every analysis except that of Fig.~5 and Extended Data Figs.~6 and~7
(Section~\ref{F-sec:fig7} of the figure methods): induction reads the true commensal fraction $C_{\mathrm{rat}}$
and reversion occurs at rate $\tau_{t,d}$. The two stability and specificity parameters stored in the parameter grids
(\texttt{phi\_instability}, \texttt{treg\_eff}) are 0 in every row and are removed
before the sweep's right-hand side is built, which contains neither; the fixed-point
finder uses the general right-hand side of \texttt{lib/ODE\_SYSTEM\_v6.R} with
$\phi = 0$ and $\eta = 1$ fixed, which is the same function.
%% from: census [grid]

%% ---- from docs/methods_parts/core.tex, lines 82-105
\subsection{Fixed constants}

\begin{table}[H]
\centering
\caption{Fixed constants (\texttt{STRUCTURAL} and \texttt{EFFERO} in
\texttt{pipeline/01\_sweep.R}). Bacterial quantities are in state units of
$10^{11}$ organisms.}
\label{tab:core-const}
\begin{tabular}{lll}
\toprule
Symbol & Value & Meaning\\\midrule
$K_E$ & 100 & epithelial capacity (cells, arbitrary unit)\\
$K_M$, $K_T$ & 100, 100 & macrophage and Treg pool sizes\\
$K_R$ & 1 & ROS scale (enters $K_{E_a}$ only)\\
$K_{L,P}$ & 1000 ($10^{14}$ organisms) & luminal pathogen carrying capacity\\
$K_{LP,P}$ & 100 ($10^{13}$ organisms) & LP pathogen carrying capacity\\
$K_{LP,C}$ & 100 ($10^{13}$ organisms) & LP commensal normalizing capacity\\
$C_L$ & 1000 ($10^{14}$ organisms) & luminal commensal reservoir\\
$a_0$ & $4.5\times10^{-3}$ & maximal efferocytosis rate per macrophage\\
$a_{\mathrm{floor}}$ & $2\times10^{-4}$ & efferocytosis rate at saturating load\\
$n_g$ & 4 & Hill exponent of $g$\\
\bottomrule
\end{tabular}
\end{table}

\subsection{Sampled kinetic parameters}
\label{sec:kin-table}

Table~\ref{tab:core-kin} lists the 22 kinetic parameters that are sampled (Section~\ref{sec:sampling}), with their symbols in the equations above and their names in the code. Three further parameters, $K_g$, $\bros$ and $r_P$, are set on grids (Section~\ref{sec:axes}).

%% ---- from docs/methods_parts/core.tex, lines 139-170
\begin{table}[H]
\centering
\caption{Sampled kinetic parameters (rates per unit time; $w$ dimensionless), their names in the code (\texttt{lib/ODE\_SYSTEM\_v6.R}, \texttt{lib/skeleton\_defs.R}), ranges and sampling scale.}
\label{tab:core-kin}
\resizebox{\linewidth}{!}{%
\begin{tabular}{lllll}
\toprule
Parameter & Code & Range & Scale & Role\\\midrule
$\kappa_E$ & \texttt{kappa\_E} & 0.05--0.15 & log & epithelial regeneration\\
$\tau_r$ & \texttt{tau\_r} & 0.05--0.50 & log & epithelial damage by ROS\\
$\tau_{P,B}$ & \texttt{tau\_P\_B} & 0.05--0.50 & log & epithelial damage by LP pathogens (necrosis)\\
$\mu_{E_a}$ & \texttt{mu\_Ea} & 0.05--0.95 & log & secondary necrosis of apoptotic cells\\
$\mu_R$ & \texttt{mu\_R} & 0.05--0.50 & log & ROS decay\\
$\tau_{01}$ & \texttt{tau\_01} & 0.005--0.95 & log & $D$-driven M1 activation\\
$\tau_{10}$ & \texttt{tau\_10} & 0.005--0.95 & log & $S$-driven M1 deactivation\\
$\tau_{1,d}$ & \texttt{tau\_1\_d} & 0.005--0.95 & log & spontaneous M1 deactivation\\
$\tau_{\mathrm{eff}}$ & \texttt{tau\_eff} & 0.5--95 & log & efferocytosis-driven M1 deactivation\\
$\tau_{t,a}$ & \texttt{tau\_t\_a} & $5\times10^{-5}$--$9.5\times10^{-3}$ & log & Treg induction\\
$\tau_{t,d}$ & \texttt{tau\_t\_d} & 0.005--0.95 & log & Treg reversion\\
$\alpha_0$ & \texttt{alpha\_0} & $5\times10^{-5}$--$10^{-4}$ & log & microbial clearance by $M_0$\\
$\alpha_{1,e}$ & \texttt{alpha\_1\_e} & $10^{-4}$--$5\times10^{-4}$ & log & microbial engulfment by $M_1$\\
$\alpha_{1,r}$ & \texttt{alpha\_1\_r} & 0.05--0.95 & log & ROS-mediated microbial clearance\\
$\alpha_{\mathrm{AMP}}$ & \texttt{alpha\_AMP} & 0.01--0.10 & log & AMP killing of luminal pathogens\\
$\gamma_{B,P}$ & \texttt{gamma\_B\_P} & 0.001--0.10 & log & pathogen translocation\\
$\gamma_{B,C}$ & \texttt{gamma\_B\_C} & 0.001--0.10 & log & commensal translocation\\
$\sigma_D$ & \texttt{sigma\_D} & 0.05--0.95 & log & danger-signal production\\
$\sigma_S$ & \texttt{sigma\_S} & $5\times10^{-4}$--$9.5\times10^{-3}$ & log & suppressive-signal production\\
$\tau_D$ & \texttt{tau\_D} & 0.001--0.50 & log & danger-signal decay\\
$\tau_S$ & \texttt{tau\_S} & 0.001--0.50 & log & suppressive-signal decay\\
$w$ & \texttt{w\_E\_over\_M} & 0.05--0.95 & linear & weight of $f_{E_n}$ against $f_{PC}$ in $dD/dt$\\
\bottomrule
\end{tabular}}
\end{table}

%% ---- from docs/methods_parts/core.tex, lines 107-119
\subsection{Units and the pathogen extinction floor}
\label{sec:units}

One state unit of bacteria is
$10^{11}$ organisms in every compartment, so one organism is $10^{-11}$ units, and a
compartment holding less than one organism holds none. Every integration therefore
carries a root on $\max(P_L, P_{LP}) - 10^{-11}$ (\texttt{EXT\_TOL}); when it is
reached, an event sets both $P_L$ and $P_{LP}$ to exactly zero
(\texttt{lib/extinction\_floor.R}, sourced by the sweep and by every figure script).
Every term of $dP_L/dt$ carries a factor $P_L$, and every term of $dP_{LP}/dt$ a factor
$P_L$ or $P_{LP}$, so both remain zero afterwards. The root is on the larger of the two
compartments because zeroing $P_{LP}$ alone while $P_L$ is above the floor leaves the
translocation influx positive. $C_{LP}$ is not floored: its influx is proportional to
the fixed reservoir $C_L$. Without the floor, a pathogen driven to $10^{-100}$ can
regrow once $M_1$ is suppressed, and the outcome is then decided by the sign of
round-off error.

\subsection{The total burden}
\label{sec:burden}

Outcomes are ordered throughout by the total burden of a state,
\begin{equation}
b(E_h, P_{LP}) = \sqrt{(K_E - E_h)^2 + P_{LP}^2},
\label{eq:burden}
\end{equation}
the Euclidean distance in the $(E_h, P_{LP})$ plane from the healthy, pathogen-free
state $(K_E, 0)$. $K_E = K_{LP,P} = 100$ puts epithelial deficit and LP pathogen on one
scale; on the sterile branch ($P_{LP} \equiv 0$) the total burden reduces to the
epithelial deficit $K_E - E_h$. It ranges from 0 to $100\sqrt2 \approx 141$. Where it is
evaluated on a trajectory rather than on a fixed point, negative round-off values of
$P_{LP}$ are replaced by 0. Its role in attractor ordering, in the gate and in the Treg
effect is defined in Sections~\ref{sec:ic-line}, \ref{sec:gate} and \ref{sec:classes}.

\subsection{Modifications used by single experiments}
\label{sec:model-mods}

Three experiments change the right-hand side above, and only there: the efferocytosis
knockout of Fig.~2 multiplies $g(L)$ by a factor $\varepsilon_{\mathrm{KO}} \in \{0, 1\}$
(Section~\ref{F-sec:fig1-ko} of the figure methods); the Treg stability and specificity sweep of Fig.~5 and
Extended Data Figs.~6 and~7 makes Treg reversion depend on $D$ and Treg induction read a noisy
commensal fraction (Section~\ref{F-sec:fig7} of the figure methods); and the treatments of Supplementary Fig.~3 delay the onset of Treg induction or rescale $\tau_{10}$, $\tau_{t,d}$
or $\tau_{t,a}$ (Section~\ref{F-sec:figS07} of the figure methods). Every other analysis integrates the
equations exactly as written above.

%% ===========================================================================
\section{Numerical methods}
\label{sec:numerics}

%% ---- from docs/methods_parts/core.tex, lines 194-213
\subsection{Numerical integration}
\label{sec:solver}

Every integration uses \texttt{deSolve::ode} with method \texttt{lsodar} (root-finding
is required by the extinction event), absolute tolerance $10^{-14}$ and relative tolerance
$10^{-8}$. The sweep uses a maximum internal step $h_{\max} = 5$ and at most $2\times10^{5}$
steps, to the horizon $T_{\max} = 12{,}000$; the settings of the other analyses are in
Table~\ref{tab:solver}. The absolute tolerance must lie at least two decades
below the extinction threshold $10^{-11}$, or the root would fire on integration error;
\texttt{lib/extinction\_floor.R} refuses to load otherwise. The maximum step is set
explicitly because \texttt{deSolve} otherwise takes it from the output spacing, which
for the two-point output grid $\{0, T_{\max}\}$ of the sweep is the whole interval.
With that default the solver stepped over the stiff tail and returned non-finite
states, $E_h > K_E$, or a different basin; on test cases $h_{\max} = 10$ still failed
where $h_{\max} = 5$ did not. The horizon was chosen from the slowest decay rates of
stable fixed points ($|\mathrm{Re}\,\lambda|$ down to about $5\times10^{-4}$): a mode
at that rate has decayed to a relative residual of about $2\times10^{-3}$ by
$t = 12{,}000$, below the $10^{-2}$ tolerance of the convergence check
(Section~\ref{sec:trust}). In the sweep an integration counts as failed if the solver
raises an error or a warning, or if its last output time is not $T_{\max}$; 73 Tregs-OFF
and 29 Tregs-ON lines among the retained combinations contain at least one failed probe
integration.
%% from: census [seed]

\subsection{Solver settings by analysis}
\label{sec:solver-table}

All analyses use the same method, tolerances, extinction floor and non-negativity
clamp; they differ in horizon, output grid, maximum step and step budget
(Table~\ref{tab:solver}). In \texttt{deSolve} the step budget (\texttt{maxsteps})
counts internal steps between two consecutive output times, so with the two-point
output grid of the sweep it bounds the whole integration. A horizon other than
$T_{\max} = 12{,}000$ is stated wherever it is used.

\begin{table}[H]
\centering
\caption{Solver settings per analysis. Common to all: \texttt{deSolve::ode}, method
\texttt{lsodar}, absolute tolerance $10^{-14}$, relative tolerance $10^{-8}$, the
pathogen extinction root at $10^{-11}$ (Section~\ref{sec:units}) and evaluation of the
right-hand side at $\max(y,0)$; the one exception is the forward stability test of a
fixed point, which uses method \texttt{lsoda} without the extinction root.
$h_{\max}$: maximum internal step; steps: maximum internal steps per output interval.}
\label{tab:solver}
\footnotesize
\begin{tabularx}{\linewidth}{@{}L l L r r l@{}}
\toprule
Analysis & Horizon & Output times & $h_{\max}$ & steps & Section\\\midrule
main sweep, bisection midpoints & 12{,}000 & $\{0, 12{,}000\}$ & 5 & $2\times10^{5}$ & \ref{sec:solver}\\
forward stability test of a fixed point (\texttt{lsoda}, no extinction root) & 20{,}000 & $\{0, 20{,}000\}$ & 5 & $10^{5}$ & \ref{sec:fp}\\
efferocytosis knockout (Figs.~2C--2G) & 12{,}000 & final state only & 5 & $10^{5}$ & \ref{F-sec:fig1}\\
Figs.~2A, 2B (example) & 4{,}000 & every 2 & 5 & $10^{5}$ & \ref{F-sec:fig1-AB}\\
load balance (Fig.~3C, Extended Data Figs.~2C and~3) & 12{,}000 & every 0.5 & 5 & $4\times10^{5}$ & \ref{F-sec:fig4}\\
basin labels at $t = 0$ (Supplementary Figs.~1B, 1C) & 12{,}000 & $\{0, 12{,}000\}$ & 5 & $10^{5}$ & \ref{F-sec:fig5}\\
Supplementary Fig.~1A (example) & 1{,}000 / 12{,}000 & 0 and 1{,}200 log-spaced from $10^{-3}$ & 5 & $10^{5}$ & \ref{F-sec:fig5}\\
damage-matched comparison (Fig.~4) & 12{,}000 & 0 and 400 log-spaced from 0.05 & 5 & $10^{5}$ & \ref{F-sec:fig6}\\
Figs.~4B, 4C (1{,}000 pairs) & 12{,}000 & 0 and 2{,}000 log-spaced from 0.05 & 5 & $10^{5}$ & \ref{F-sec:fig6}\\
Fig.~4D (one unit, four arms) & 12{,}000 & 0 and 3{,}000 log-spaced from 0.05 & 5 & $10^{5}$ & \ref{F-sec:fig6}\\
first peak of each arm (Fig.~4, Extended Data Fig.~5) & 12{,}000 & 0 and 2{,}000 log-spaced from 0.05 & 5 & $10^{5}$ & \ref{F-sec:first_peak}\\
Treg stability and specificity (Fig.~5, Extended Data Figs.~6 and~7) & 12{,}000 & every 1 & 1 & $2\times10^{5}$ & \ref{F-sec:fig7}\\
phase portraits (Figs.~3A, 3B; Extended Data Figs.~2A, 2B; Supplementary Figs.~2A, 2B) & 12{,}000 & every 2 & 5 & $10^{5}$ & \ref{F-sec:figS01}\\
Treg treatments (Supplementary Figs.~3B, 3C) & 25{,}000 & 0 and 300 log-spaced from 0.05 & 5 & $10^{5}$ & \ref{F-sec:figS07}\\
Supplementary Fig.~3A & 12{,}000 & 300 linear on $[0,300]$, 120 log-spaced on $[300, 12{,}000]$ & 5 & $10^{5}$ & \ref{F-sec:figS07}\\
Supplementary Figs.~3D--3G & 25{,}000 & 0 and 400 log-spaced from 0.05 & 5 & $10^{5}$ & \ref{F-sec:figS07}\\
\bottomrule
\end{tabularx}
\end{table}

%% ---- from docs/methods_parts/core.tex, lines 215-257
\subsection{Fixed points and their stability}
\label{sec:fp}

Fixed points are computed separately for each Treg condition by the reductions in
\texttt{lib/FP\_finder\_v6.R}.

\emph{Sterile branch} ($P_L = P_{LP} = 0$). At a fixed point every state is a closed-form
function of the dead-cell load $L$: the damage drive is $A = K_M\,L\,g(L)$,
$M_1 = A\kappa_E / \big(a\,(\kappa_E K_E - A)\big)$ with $a = \tau_r\beta_r/\mu_R$,
and $R$, $E_h$, $E_a$, $E_n$, $C_{LP}$, $D$ (and, with $\theta = 1$, $T_r$ and $S$) follow.
Fixed points are the roots of the scalar residual $G(L) = dM_1/dt$ evaluated on that
state. $G$ is evaluated on a uniform grid on
$[0,\ 1.2\,\kappa_E K_E/(K_M a_{\mathrm{floor}})]$ with $\max(6000, 12\,L_{\max})$ points,
plus a logarithmic sub-scan of the first grid cell (400 points from $10^{-10}$), with
refinement where $G$ becomes undefined; every sign change is solved by
\texttt{uniroot} and accepted if $|G| \le 10^{-6}$. $L = 0$ is an exact root (the healthy
state $E_h = K_E$). Supplied seeds (below) are checked for roots the grid did not bracket.

\emph{Pathogenic branch.} The fixed-point conditions reduce to three residuals in
$(M_1, P_{LP}, L)$ on each of the two branches of $dP_L/dt = 0$ ($P_L = 0$, or
$P_L = K_{L,P}(1 - \alpha_{\mathrm{AMP}} D / r_P)$), solved by damped Newton iteration
(central-difference Jacobian, backtracking line search, convergence at $10^{-9}$,
acceptance at $10^{-6}$). In the grid-seeded mode Newton starts from
$M_1 \in \{1, 32.3, 63.7, 95\}$ $\times$ $P_{LP} \in \{0, 0.2, 0.5, 1, 48, 95\}$ $\times$
$L \in \{1, 5, 30, 100, 400\}$ on both $P_L$ branches, then from every solution found with
$P_{LP}$ raised to at least 0.2, 0.5 and 2 on the $P_L > 0$ branch. Supplied seeds are
added. Solutions are de-duplicated after rounding $(E_h, P_{LP}, M_1)$ to three decimals.

\emph{Endpoint seeding.} Both finders accept the final states of integrated
trajectories as seeds (columns $M_1$, $P_{LP}$, $L$). A final state lies on the attractor
it reached, so such seeds can only add fixed points.

\emph{Stability.} Each fixed point is classified from the eigenvalues of a reduced
Jacobian (central differences) in which $M_0 = K_M - M_1$ and $T_0 = K_T - T_r$ are
eliminated, removing the two zero eigenvalues of the conservation laws; on the sterile
branch the two pathogen directions are dropped as well (9 variables; 11 on the
pathogenic branch). On the sterile
branch the fixed point is \emph{stable} if the largest real part is below $-10^{-8}$, a
\emph{saddle} if it is above $10^{-8}$, and \emph{marginal} otherwise. On the pathogenic
branch it is stable below $-10^{-8}$ and a saddle above $10^{-4}$; in between, a forward
test decides: the fixed point is perturbed by $10^{-3}(|x|+1)$ in every state,
integrated to $t = 20{,}000$ with the tolerances and maximum step of
Section~\ref{sec:solver} (method \texttt{lsoda}, at most $10^{5}$ steps, no extinction
root; Table~\ref{tab:solver}), and called
stable if the scaled distance at the end is at most five times the initial one (saddle
if larger, marginal if the integration fails). Only fixed points labeled stable are
attractors anywhere below.

%% ---- from docs/methods_parts/core.tex, lines 259-351
\subsection{The initial-condition line, basin assignment and thresholds}
\label{sec:ic-line}

\subsubsection{Initial conditions}

Every sweep integration starts on one line of initial conditions indexed by the initial
healthy epithelium $e = E_h(0)$:
\begin{equation}
E_h = e,\quad E_n = K_E - e,\quad E_a = 0,\quad M_0 = K_M,\quad T_0 = K_T,\quad
M_1 = T_r = R = D = S = C_{LP} = 0,
\label{eq:icline}
\end{equation}
and on the pathogenic branch $P_L = 500$, $P_{LP} = 10$ (both zero on the sterile
branch). The injury is necrotic debris equal to the missing epithelium. $e = 100$ is
excluded: in the sterile system it is itself a fixed point (a saddle whose unstable
direction lies along $C_{LP}$), so a trajectory starting there belongs to no basin;
$e = 99$ leaves its neighborhood long before $T_{\max}$.

\subsubsection{Probe, fixed-point set and labels}

For each combination and condition the line is probed at
$e \in \{0, 10, 20, \dots, 90, 99\}$ (11 values), each trajectory integrated to $T_{\max}$.
The \emph{grid-seeded} fixed-point set (sterile: the full $L$ scan; pathogenic: the
grid-seeded Newton) is computed once per condition and records what exists, saddles
included. The \emph{endpoint-seeded} set is then computed with the 11 final states as
seeds; on the pathogenic branch it uses no grid seeds, and if it is empty the
grid-seeded set is used instead. Among the 2,248,846 retained pathogenic combinations
this fallback was taken on 56,609 Tregs-OFF and 69,389 Tregs-ON lines; it was never
needed on the sterile branch.
%% from: census [seed]

Each final state is assigned to the nearest fixed point of the endpoint-seeded set in
the scaled squared distance $\sum_v \big((x_v - x_v^*)/(|x_v^*| + 1)\big)^2$ over
$v \in \{E_h, E_a, E_n, P_{LP}, C_{LP}, M_1, R, D, T_r, S, P_L\}$. The assignment counts
only if that fixed point is stable; a final state matched to a saddle, a marginal
point or nothing is \emph{unassigned}.

\subsubsection{Attractor clustering}
\label{sec:clustering}

\emph{Attractor identity.} Two stable fixed points are the same attractor when they
agree to within 1 in both $E_h$ and $P_{LP}$ (a Chebyshev ball of radius 1 on
$(E_h, P_{LP})$; \texttt{lib/attractor\_cluster.R}), grouped by single linkage
(union--find), which is independent of input order. $K_E = K_{LP,P} = 100$ puts the two
coordinates on one scale, and one unit of $P_{LP}$ is $10^{11}$ organisms, the scale at
which two pathogen loads are no longer the same outcome. There is no separate
extinction clause: a pathogen-free fixed point and one carrying $P_{LP} \le 1$ fall in
one attractor if they also agree in $E_h$. Attractors are ordered by the total burden $b$ (Eq.~\ref{eq:burden}); the \emph{low-burden attractor} is the group containing the lowest-burden reached
fixed point, and each group is represented by its lowest-burden member. The
total burden orders outcomes; it is not used to decide identity, because two states that
trade epithelial damage against pathogen load can share a burden value (measured
case: $E_h = 38.7, P_{LP} = 0$ against $E_h = 43.7, P_{LP} = 26.9$, burdens 61.30 and
62.40).

\emph{Labels.} Each probed $e$ is labeled 1 if its final state belongs to the low-burden attractor,
0 if it belongs to any other attractor, and is left unlabelled otherwise. The number
of label changes along the labeled values, $n_{\mathrm{cross}}$, and the number of
distinct attractors reached, $n_{\mathrm{ray}}$, are recorded; $n_{\mathrm{ray}} = 1
\Leftrightarrow n_{\mathrm{cross}} = 0$ holds in every retained combination and
condition. At least one probed value was unassigned on 1,528 Tregs-OFF
%% from: census [seed]
and 3,633 Tregs-ON lines.

\subsubsection{\texorpdfstring{The $E_h(0)$ threshold by bisection}{The Eh(0) threshold by bisection}}
\label{sec:bisection}

When the labels change exactly once, the bracketing pair of probe values is bisected
(each midpoint integrated, assigned and labeled by the same rule) until the bracket is
at most \texttt{EH0\_TOL} $= 0.1$ wide, and the threshold $E_h(0)^*$ is the midpoint of
the final bracket. A bracket of width 10 needs seven halvings, so thresholds lie on a
lattice of step $10/2^7 = 0.078125$ ($9/2^7 = 0.0703$ in the bracket $[90, 99]$). The side
on which the low-burden attractor lies is read off the labels and recorded as the
\emph{direction}: \emph{increasing} if the low-burden attractor is reached above the
threshold. If the labels change more than once no threshold is reported
(\texttt{non\_monotonic\_}$k$); if a bisection midpoint cannot be labeled, none is
reported either (\texttt{bisect\_unlabelled}). Tregs-OFF outcomes among the retained
combinations: sterile 215,192 single attractor, 160,200 one crossing, 21 more than one
crossing, 3 unlabelled midpoint; pathogenic 1,172,831, 1,037,904, 36,956 and 1,155.
%% from: census [seed]

\subsubsection{Convergence check and exclusion}
\label{sec:trust}

For each probed $e$ with a stable assignment, the relative deviation of the final state
from its assigned fixed point is $\max_v |x_v - x^*_v|/(|x^*_v|+1)$ over all 13 states.
A combination is retained only if the median of this deviation over its probed values
is at most $10^{-2}$ in both conditions (\texttt{pipeline/02\_assemble.R}). This
excluded 17,769 of 2,642,031 combinations (0.67\%; Tregs-OFF median above the
tolerance in 11,097, Tregs-ON in 14,402, no stable assignment at all in 150 and 1,305),
leaving 2,624,262: 375,416 sterile and 2,248,846 pathogenic. The excluded fraction is
0.44--1.11\% per $r_P$ level (highest at $r_P = 0.08$), 0.34--1.06\% per $K_g$ level
(highest at 25--30) and 0.63--0.71\% per $\beta_r$ level.
%% from: census [trust]
At assembly, fixed points of one combination within the attractor-identity tolerance
are merged (state vector averaged; the merged point is stable only if every member is)
and the probe map is re-indexed. In the stored attractor coordinates, values of $P_{LP}$
in $[-10^{-9}, 0)$ are set to 0 and values below $-10^{-9}$ to missing
(\texttt{lib/plp\_clean.R}).

%% ===========================================================================
\section{Parameter sampling and the sweep}
\label{sec:sampling}

%% ---- from docs/methods_parts/core.tex, lines 124-137
\subsection{Kinetic parameters}

Twenty-two kinetic parameters were drawn by Latin hypercube sampling in five grids of
2,000 vectors each (\texttt{pipeline/00\_parameter\_grids.R}). For each parameter, the
unit interval is divided into 2,000 equal strata, one uniform draw is taken in each,
and the strata are randomly permuted, independently per parameter. Each strictly
positive rate constant is mapped log-uniformly onto its range
($\exp(\log a + u\,(\log b - \log a))$); the dimensionless weight $w$ is mapped
linearly. Table~\ref{tab:core-kin} gives the ranges; the observed minima and maxima of
the 10,000 stored vectors lie inside them. The generating script set
%% from: census [grid]
no random seed; the five stored grids (\texttt{data/tau\_grid\_v6\_1.rds} to
\texttt{\_5.rds}) are the only record of the parameter sets and every stage reads them.
Each vector carries a globally unique key $\texttt{tau\_id} = 10^4 \times
\text{grid} + \text{row}$, stamped into the grid when it is drawn.

%% ---- from docs/methods_parts/core.tex, lines 172-192
\subsection{Swept axes and parameter combinations}
\label{sec:axes}

Three parameters are set on a grid applied to every kinetic vector
(\texttt{pipeline/01\_sweep.R}): the efferocytosis half-saturation load
$K_g \in \{15, 20, 25, 30, 35, 40, 45\}$; the ROS production rate
$\beta_r \in \{0.0015, 0.0025, \dots, 0.0095\}$ (nine levels, step 0.001);
and the pathogen growth rate $r_P \in \{0, 0.08, 0.16, 0.24, 0.32, 0.40, 0.48\}$. A
\emph{parameter combination} is a (kinetic vector, $K_g$, $\beta_r$, $r_P$)
quadruple, 441 per kinetic vector, keyed by
(\texttt{tau\_id}, $K_g$, $\beta_r$, $r_P$).

\emph{Branches.} $r_P = 0$ defines the \emph{sterile} branch: both pathogen
compartments start at zero and remain there. $r_P > 0$ defines the
\emph{pathogenic} branch, which is divided after the gate (Section~\ref{sec:classes})
into \emph{pathogenic help} and \emph{pathogenic harm} by the effect of Tregs on the
low-burden attractor.

\emph{Coverage.} The sweep was run on grids 1--3 (6,000 kinetic vectors, $6{,}000 \times 441 = 2{,}646{,}000$ scheduled combinations; grids 4 and 5
were not swept, but the efferocytosis knockout of Fig.~2 draws from all five grids). 5,991 vectors returned all 441 combinations and nine returned none
(grid 1 row 1885; grid 3 rows 180, 184, 185, 350, 989, 990, 1024, 1025), giving
2,642,031 of 2,646,000 combinations. No vector returned a partial set.
%% from: census [sweep]

%% ===========================================================================
\section{Classification and the gate}
\label{sec:classification}

%% ---- from docs/methods_parts/core.tex, lines 353-375
\subsection{The bistability gate}
\label{sec:gate}

Every Treg analysis is gated on the Tregs-OFF condition (criteria (1)--(3);
\texttt{pipeline/03\_select.R}, \texttt{GAP\_CUT} $= 20$) together with one
commensurability condition on the Tregs-ON line (criterion (4)). A combination passes if
\begin{enumerate}
\item the lowest- and highest-burden Tregs-OFF attractors reached from the line are more
  than 20 burden units apart, $b_{\mathrm{hi}} - b_{\mathrm{lo}} > 20$ (the total burden
  ranges from 0 to $100\sqrt2 \approx 141$); this requires two reached attractors;
\item the Tregs-OFF labels change exactly once ($n_{\mathrm{cross}} = 1$), so a single
  Tregs-OFF threshold exists;
\item the Tregs-OFF direction is increasing: more healthy epithelium at $t = 0$ reaches
  the low-burden attractor;
\item the Tregs-ON line is commensurable with it: either the Tregs-ON threshold is a
  single increasing crossing, or the Tregs-ON line reaches the low-burden attractor from
  every probed value. In the kinds of Section~\ref{sec:gap} this keeps \texttt{measured}
  and \texttt{on\_lo\_basin\_throughout} and discards \texttt{opposite\_directions} and
  \texttt{none}; the rationale and the per-class effect are given in
  Section~\ref{sec:excl_decreasing}.
\end{enumerate}
The decreasing configuration is a pathogenic phenomenon whose frequency rises with $r_P$
(36\% of the combinations passing (1)--(2) at $r_P = 0.08$, 66\% at $r_P = 0.48$);
%% from: census [gate]
it is not commensurable with the increasing one, because the resolving
set is $[E_h(0)^*, 99]$ in one and $[0, E_h(0)^*]$ in the other. The gate uses the
total burden separation; the Chebyshev separation $\max(|\Delta E_h|, |\Delta P_{LP}|)$
gives the same decision on every sterile combination with two attractors but exceeds 20
on 384,675 pathogenic combinations against 115,153 for the total burden; the total burden was kept.
%% from: census [gate]

Table~\ref{tab:gate} follows every parameter combination from the sweep design to the
classes and rescue-gap kinds of Section~\ref{sec:classes} and~\ref{sec:gap}, with the
count after every step, per branch. The Treg effect (help, neutral, harm;
Section~\ref{sec:classes}) is computed for every retained combination, so the columns
by effect are shown from the convergence check onward; the branches are formed from
the gated set only, and the effect columns of the earlier rows are for orientation.
The gated set (\texttt{data/selected\_big.rds}) was recomputed from the assembled
batches with identical keys and identical class labels. It comes from
%% from: census [check]
3{,}139 kinetic vectors (sterile 3{,}133, pathogenic 1{,}035). The pass rate rises with
$\bros$ (sterile 15.95\% at 0.0015 to 26.70\% at 0.0095; pathogenic 0.85\% to 3.55\%),
peaks at intermediate $K_g$ (sterile 34.86\% at $K_g = 25$, 4.04\% at 45; pathogenic
3.84\% at 35, 0.56\% at 15) and falls with $r_P$ (6.55\% at 0.08 to 0.76\% at 0.48).

\begin{table}[H]
\centering
\caption{The consolidated gate: parameter combinations after every step, per branch and
Treg effect. Columns: sterile (all, and by Treg effect; harm does not occur on the
sterile branch), pathogenic (all, and by Treg effect: help $=$ \emph{pathogenic help}
branch, harm $=$ \emph{pathogenic harm} branch, neutral $=$ no branch), and total. Gate
steps (1)--(3) are applied to the Tregs-OFF line only (Section~\ref{sec:gate}).
``Measured'', ``opposite'', ``none'' and ``censored'' are the rescue-gap kinds of
Section~\ref{sec:gap}. Bold: the populations of the three branches. Sources: census
[sweep], [trust], [gate], [class], [gap]; per-effect rescue-gap kinds and per-branch
returned and excluded counts from \texttt{data/selected\_big.rds} and the sweep coverage
(5{,}991 complete kinetic vectors $\times$ 63 sterile or 378 pathogenic combinations). The class rows below the gate sum to row~(3) rather than to
the bolded row~(4), because criterion~(4) is applied after the classification: gated =
measured + ON\_mono.}
\label{tab:gate}
\footnotesize
\setlength{\tabcolsep}{3.5pt}
\resizebox{\linewidth}{!}{%
\begin{tabular}{@{}lrrrrrrrr@{}}
\toprule
 & \multicolumn{3}{c}{sterile} & \multicolumn{4}{c}{pathogenic} & \\
\cmidrule(lr){2-4}\cmidrule(lr){5-8}
step & all & help & neutral & all & help & harm & neutral & total\\\midrule
scheduled (grids 1--3, 6{,}000 vectors $\times$ 441) & 378{,}000 & & & 2{,}268{,}000 & & & & 2{,}646{,}000\\
returned by the sweep (5{,}991 vectors) & 377{,}433 & & & 2{,}264{,}598 & & & & 2{,}642{,}031\\
\quad excluded by the convergence check & 2{,}017 & & & 15{,}752 & & & & 17{,}769\\
retained & 375{,}416 & 175{,}351 & 200{,}065 & 2{,}248{,}846 & 399{,}870 & 547{,}474 & 1{,}301{,}502 & 2{,}624{,}262\\
(1) $b_{\mathrm{hi}} - b_{\mathrm{lo}} > 20$ (Tregs OFF) & 89{,}330 & 43{,}377 & 45{,}953 & 115{,}153 & 21{,}710 & 33{,}399 & 60{,}044 & 204{,}483\\
(2) one Tregs-OFF crossing & 89{,}330 & 43{,}377 & 45{,}953 & 111{,}695 & 21{,}207 & 32{,}109 & 58{,}379 & 201{,}025\\
(3) Tregs-OFF direction increasing & 89{,}328 & 43{,}375 & 45{,}953 & 57{,}497 & 17{,}495 & 5{,}246 & 34{,}756 & 146{,}825\\
(4) Tregs-ON line commensurable: \textbf{gated} & \textbf{89{,}322} & 43{,}371 & 45{,}951 & \textbf{52{,}108} & \textbf{14{,}995} & \textbf{2{,}922} & 34{,}191 & \textbf{141{,}430}\\
\midrule
class bistable & 78{,}169 & 33{,}491 & 44{,}678 & 54{,}784 & 16{,}225 & 4{,}066 & 34{,}493 & 132{,}953\\
\quad gap measured & 78{,}163 & 33{,}487 & 44{,}676 & 49{,}395 & 13{,}725 & 1{,}742 & 33{,}928 & 127{,}558\\
\quad gap opposite directions & 0 & 0 & 0 & 5{,}298 & 2{,}443 & 2{,}313 & 542 & 5{,}298\\
\quad gap none & 6 & 4 & 2 & 91 & 57 & 11 & 23 & 97\\
class ON\_mono (gap censored) & 11{,}159 & 9{,}884 & 1{,}275 & 2{,}713 & 1{,}270 & 1{,}180 & 263 & 13{,}872\\
class mono\_A, unresolved\_on & 0 & 0 & 0 & 0 & 0 & 0 & 0 & 0\\
\bottomrule
\end{tabular}}
\end{table}
%% from: census [sweep] [trust] [gate] [class] [gap]; per-branch returned = 5,991 x 63 and
%% 5,991 x 378 (no vector returned a partial set); excluded = returned - retained.
%% Rescue-gap kinds by effect: selected_big.rds, dcast(rescue_gap_kind ~ branch x treg_effect),
%% recomputed at assembly (scratch chk3.R); they reproduce census [gap] for the three branches.

%% ---- from docs/methods_parts/core.tex, lines 409-413
Losses at each step: (1) removed 286,086 sterile and 2,133,693 pathogenic combinations
(one attractor reached, 215,192 and 1,172,831, or two attractors 20 or fewer burden
units apart); (2) removed 3,458 pathogenic combinations with two Tregs-OFF crossings and
no sterile one; (3) removed 2 sterile combinations (bisection midpoint unlabelled, no
direction) and 54,198 pathogenic ones (53,881 decreasing, 317 unlabelled midpoint).

%% ---- from docs/methods_parts/core.tex, lines 415-444
\subsection{What Tregs do: classes, effect and branches}
\label{sec:classes}

\emph{Classes.} Each gated combination is classified by its Tregs-ON line:
\begin{itemize}
\item \textbf{bistable}: the Tregs-ON line reaches two or more attractors;
\item \textbf{ON\_mono} (\texttt{mono\_B} in the code): the Tregs-ON line reaches a
  single attractor, and that attractor is the low-burden one. This is recognized in one
  of two ways: the sweep records the low-burden attractor as reached from every probed
  value, or, failing that, the single Tregs-ON attractor lies closer in total burden to the
  Tregs-OFF low-burden attractor than to the Tregs-OFF high-burden one;
\item \textbf{mono\_A}: a single Tregs-ON attractor that lies closer to the Tregs-OFF
  high-burden attractor;
\item \textbf{unresolved\_on}: no Tregs-ON attractor was recovered.
\end{itemize}
Every one of the 13,872 ON\_mono combinations was labeled by the first rule; mono\_A
and unresolved\_on did not occur. ON\_mono is a statement about the initial-condition
line only: a second stable Tregs-ON fixed point exists off the line (grid-seeded set)
in 763 of 11,159 sterile, 581 of 1,270 pathogenic help and 554 of 1,180 pathogenic harm
ON\_mono combinations.
%% from: census [class] [gap]

\emph{Treg effect.} The effect of Tregs is judged on the low-burden attractor:
$\Delta b_{\mathrm{lo}} = b(E_h^{\mathrm{lo,ON}}, P_{LP}^{\mathrm{lo,ON}}) -
b(E_h^{\mathrm{lo,OFF}}, P_{LP}^{\mathrm{lo,OFF}})$, recomputed from the stored attractor
coordinates. It is \emph{help} if $\Delta b_{\mathrm{lo}} < -1$, \emph{harm} if
$\Delta b_{\mathrm{lo}} > 1$, \emph{neutral} otherwise (effect tolerance $\delta = 1$,
\texttt{EFFECT\_TOL}). The
total burden weights one unit of epithelial damage equally with one unit of LP pathogen;
$K_E = K_{LP,P}$ puts the two on the same scale. The effect answers a different question
from the rescue gap below, and the two can disagree in the same combination.

\paragraph{Coverage of the help class.} The help class is not confined to a corner of the
grid. Among the gated pathogenic combinations it is populated at every pathogen growth rate,
from 35.0\% of combinations at $r_P = 0.08$ falling monotonically to 22.3\% at $r_P = 0.48$
(8{,}558 of 24{,}433 down to 633 of 2{,}842), and at every ROS production rate, between 27.2\%
at $\beta_r = 0.0015$ and 32.7\% at $\beta_r = 0.0035$ (579 of 2{,}127 and 1{,}752 of 5{,}351),
as well as at every $K_g$. Jointly the coverage is almost but not quite complete: 53 of the 54
$(\beta_r, r_P)$ cells hold help combinations, the exception being the corner of the weakest
response against the fastest pathogen, $\beta_r = 0.0015$ with $r_P = 0.48$, which holds 39
combinations and no help; its neighbor $(0.0015, 0.40)$ holds one of 67.
% from: scratchpad chk_help.R, on data/selected_big.rds

\emph{Branches.} \emph{Sterile}: all gated combinations with $r_P = 0$.
\emph{Pathogenic help} and \emph{pathogenic harm}: gated combinations with $r_P > 0$ and
that effect. The 34,191 gated pathogenic neutral combinations belong to no branch.

%% ---- from docs/methods_parts/core.tex, lines 459-462
The gated sterile branch holds 78,163 bistable and 11,159 ON\_mono combinations (help and
neutral pooled; harm does not occur); gated pathogenic help 13,725 and 1,270 (14,995);
gated pathogenic harm 1,742 and 1,180 (2,922). The branch files
\texttt{data/skeleton\_<branch>.rds} and \texttt{data/rescue\_<branch>.rds} were written before
criterion~(4) and hold the larger criterion-(3) populations: 78,169 and 11,159 sterile,
16,225 and 1,270 pathogenic help, 4,066 and 1,180 pathogenic harm.

%% ---- from docs/methods_parts/core.tex, lines 464-511
\subsection{The rescue gap}
\label{sec:gap}

For a combination whose Tregs-OFF and Tregs-ON thresholds are both located and run in
the same direction, the rescue gap is
\begin{equation}
\Delta E_h(0)^* = \begin{cases}
E_h(0)^*_{\mathrm{OFF}} - E_h(0)^*_{\mathrm{ON}}, & \text{increasing},\\
E_h(0)^*_{\mathrm{ON}} - E_h(0)^*_{\mathrm{OFF}}, & \text{decreasing},
\end{cases}
\label{eq:gap}
\end{equation}
the additional range of initial epithelial damage from which the tissue reaches the
low-burden attractor when Tregs are present (kind \texttt{measured}). Thresholds running
in opposite directions are not comparable and give no gap (kind
\texttt{opposite\_directions}); a combination with a missing Tregs-ON threshold and a
Tregs-ON line that is not entirely low-burden has none either (kind \texttt{none}). When
the Tregs-ON line reaches the low-burden attractor from every probed value (all ON\_mono
combinations), the recorded value is the whole resolving extension of the line,
$E_h(0)^*_{\mathrm{OFF}} - 0$, flagged as censored (kind \texttt{on\_lo\_basin\_throughout}).
Because both thresholds lie on the bisection lattice, measured gaps are multiples of
0.078125 (one sterile gap of 78,163 lies on the $[90,99]$ lattice instead), and a gap of
0 means a shift smaller than one lattice step, not the absence of a shift.

\begin{table}[H]
\centering
\caption{Rescue gaps of the bistable class per branch (percentages of the
measured gaps). All measured gaps run increasing in both arms.}
\label{tab:core-gap}
\resizebox{\linewidth}{!}{%
\begin{tabular}{lrrr}
\toprule
 & sterile (help + neutral) & pathogenic help & pathogenic harm\\\midrule
bistable & 78,169 & 16,225 & 4,066\\
\quad measured & 78,163 & 13,725 & 1,742\\
\quad opposite directions & 0 & 2,443 & 2,313\\
\quad none & 6 & 57 & 11\\
measured: exactly 0 & 53.1\% & 58.7\% & 65.7\%\\
measured: $\le$ one step (0.078125) & 73.3\% & 78.1\% & 81.6\%\\
measured: $> 0$ & 46.9\% & 41.3\% & 34.0\%\\
measured: $< 0$ (count) & 0 & 3 & 4\\
90th / 99th percentile & 0.55 / 4.06 & 0.31 / 2.42 & 0.31 / 3.14\\
maximum & 48.59 & 33.28 & 17.11\\
ON\_mono (censored) & 11,159 & 1,270 & 1,180\\
\quad recorded extension, median (range) & 27.54 (0.04--83.95) & 25.78 (0.04--98.61) & 28.16 (0.04--80.98)\\
\bottomrule
\end{tabular}}
\end{table}
%% from: census [gap] and [skel]; 90th/99th percentiles printed as 0.5469/4.0625,
%% 0.3125/2.4219, 0.3125/3.1391.

%% ---- from docs/methods_parts/fig02_03.tex, lines 464-469
\emph{The rescue-gap floor.} A measured gap is said to be \emph{at the floor} if
$\gap \le 1.01\times0.078125$, which includes the exact zeros; in the bisection such a pair
of thresholds was not separated. The floor fractions of Table~\ref{tab:core-gap} (exactly 0:
53.1\% sterile, 58.7\% pathogenic help; at most one lattice step: 73.3\% and 78.1\%), the
three negative pathogenic help gaps and the maxima 48.59 and 33.28 were obtained
independently from the branch files by \texttt{pipeline/24\_numbers\_load\_balance\_boundary.R} and
agree. Supplementary Fig.~1 excludes gaps at the floor (Section~\ref{F-sec:fig5} of the figure methods).
%% from: pipeline/24_numbers_fig04_05.R, block "gap_floor"; census [gap]; re-checked at assembly
%% from selected_big.rds (exact zeros 41,477/78,163, 8,059/13,725, 1,145/1,742)."

%% ===========================================================================
\subsection{Exclusion of combinations whose Tregs-ON threshold inverts}
\label{sec:excl_decreasing}

A combination is retained for analysis only if its Tregs-ON threshold is a single
increasing crossing, or if the Tregs-ON line reaches the low-burden attractor from every
probed value. In the kinds of Section~\ref{sec:gap} this keeps \texttt{measured} and
\texttt{on\_lo\_basin\_throughout} and discards \texttt{opposite\_directions} and
\texttt{none}.

The reason is interpretive rather than numerical. In the discarded
\texttt{opposite\_directions} combinations the Tregs-OFF threshold runs increasing, so the
low-burden attractor is reached from the healthier initial conditions as expected, while
the Tregs-ON threshold runs decreasing, so with Tregs the low-burden attractor is reached
from the \emph{more} injured ones. A regime in which a larger initial injury yields the
better outcome is outside the range of behaviour we interpret, and we remove it rather
than report it. The runs themselves are clean: each has exactly one crossing per arm and
no solver note, and every attractor coordinate is finite, so this is a modelling
judgement and not a repair of failed integrations. The \texttt{none} kind is discarded on
the separate ground that its Tregs-ON arm could not be labelled at all.

The criterion applies to every branch and to both Treg-effect classes, so that no
comparison is made between differently filtered populations. Its effect is confined to
the pathogenic branch: the sterile branch contains no combination with a decreasing
Tregs-ON threshold and loses six of 89,328 to the \texttt{none} kind. In the pathogenic
branch it removes 2,313 of 4,066 bistable harm combinations and 2,443 of 16,225 bistable
help combinations, together with 542 neutral and the 91 \texttt{none} combinations across
the three classes, leaving 2,922 harm, 14,995 help and 34,191 neutral combinations of
52,108 in total. ON\_mono combinations are unaffected, since their recorded extension is
the whole line rather than a difference between two crossings.

The damage-matched experiment of Figure~4 inherits the same criterion, which reduces its
sample from 11,940 pairs to 8,789 and the median $\beta_r^{*}/\beta_r$ from 2.19 to 1.89.
The excluded pairs are those for which the largest increase in aggression is affordable:
taken alone their median ratio is 5.13. The qualitative results are unchanged.
$\beta_r^{*} > \beta_r$ in every retained pair, the first pathogen peak is lower with
Tregs in every retained pair, and the cumulative load over the matching window is lower
with Tregs in 89.3\% of them, against 86.5\% before the exclusion.

The criterion was already in force wherever a rescue gap is plotted or averaged, since
those panels require a finite gap (\texttt{lib/fig\_skeleton.R}); stating it here makes
the counts reported in the text agree with the sample sizes shown in the figures. It is
\emph{not} a statement about how often such a regime arises, for the same reason given
for every other count in this work.
%% from: lib/fig_damage_matched.R (EXCLUDE_DECREASING, default on); counts re-derived
%% from selected_big.rds and the skeleton branch files 2026-11-06; the damage-matched
%% figures were re-rendered and the unfiltered run reproduces every pre-exclusion number.

%% ===========================================================================
\section{Statistical analysis}
\label{sec:stats}

All results are properties of deterministic (or, in Fig.~5, stochastic) simulations of
the model over sampled parameter sets. The sampled parameter space has no known
population distribution, so every proportion is a proportion of the sample or of the
stated population of parameter combinations or units, not a frequency in a biological
population, and no inferential statement refers to a population beyond the sample. The
statistics used, and where:

\begin{itemize}
\item \emph{Medians, interquartile ranges and ranges} summarize every distribution in
  the endpoint tables. Quantiles are R \texttt{quantile} type 7 unless stated.
\item \emph{Percentile bootstrap} (Figs.~2F, 2G; Section~\ref{F-sec:fig1-FG} of the figure methods): 95\%
  interval of the subsampled median $S(x)$ from 400 subsamples of 13 units per stratum
  drawn with replacement (units, not pairs), seed 1.
\item \emph{Spearman rank correlation}: Supplementary Figs.~1B, 1C, $G_X$ against the rescue gap
  (\texttt{cor.test}, \texttt{exact = FALSE}, asymptotic $p$; Table~\ref{F-tab:gainstats} of the figure methods);
  Supplementary Fig.~3B, $d^\ast$ against $t_{\mathrm{dip}}$ ($\rho = 0.994$),
  Table~\ref{F-tab:harm_dose} of the figure methods, least sufficient ladder step against $t_{\mathrm{dip}}$,
  and the diagnostic correlations of Section~\ref{F-sec:figS07} of the figure methods (coefficient only,
  R \texttt{cor}; no $p$ value is computed or reported).
\item \emph{Ordinary least squares on log--log axes}: Supplementary Figs.~1B, 1C
  ($\log_{10}\gap$ on $\log_{10} G_X$, drawn lines, no interval); Supplementary Fig.~3B (slope of $\log_{10} d^\ast$ on $\log_{10} t_{\mathrm{dip}}$ with its 95\%
  confidence interval, 1.126, 1.107--1.146).
\item \emph{Two-sample Kolmogorov--Smirnov test} (Section~\ref{F-sec:fig5} of the figure methods): the rescue
  gaps of the bistable combinations in the measured sterile subsample against the whole
  sterile pool ($D = 0.0047$, $p = 0.95$), used to check that the subsample is random,
  not to test a hypothesis about the model.
\item \emph{Wilson score intervals} (95\%; the two figures of the appendix, panels U--X)
  for the fraction of combinations not harmed at each treatment level, treating the
  combinations as independent draws; all levels are evaluated on the same combinations,
  so the intervals describe each level separately and are not a test between levels.
\item \emph{Kernel density estimates}: bivariate Gaussian (\texttt{MASS::kde2d}, default
  bandwidth) for the highest-density regions and modes of Figs.~2C--2E; univariate
  Gaussian (R \texttt{density}, default bandwidth) for the violins of Extended Data Fig.~1
and Supplementary Fig.~2. The histograms of Extended Data Fig.~5 are counts in
  bins on a pseudo-logarithmic axis, not densities.
\item \emph{Quantile fans} (Figs.~5A, 5B; Extended Data Figs.~6A, 6B; Extended Data Figs.~7C--7F): nine nested
  central quantile intervals across units.
\end{itemize}

No other test is performed. Fig.~3C, Extended Data Figs.~2C and~3, Fig.~4, Fig.~5,
Extended Data Fig.~6 and Extended Data Fig.~7 contain no test or interval other than those listed; in
particular Fig.~4 contains no bootstrap, test or correlation, and Fig.~3C and Extended Data Fig.~2C shows a
random display subset of 20 trajectories per branch on which no statistic is computed.
No correction for multiple comparisons is applied, because no family of hypothesis
tests is interpreted.

\emph{Random seeds.} Unit and combination draws: Fig.~2, seed 1
(\texttt{pipeline/06\_effero\_knockout.R}); Fig.~3C, Extended Data Figs.~2C and~3,
seed 1 (\texttt{pipeline/20\_load\_balance.R}), display subset seed 1 re-set per branch;
Supplementary Fig.~1, pool permutation seed 20260820 (\texttt{pipeline/21\_corridor\_tzero.R}),
display subsample seed 1; Fig.~4, unit draw seed 1
(\texttt{pipeline/43\_damage\_matched\_gated\_units.R}), trajectory subset seed 1
(\texttt{pipeline/41\_damage\_matched\_trajectories.R}), label placement seed 1
(\texttt{ggrepel});
Fig.~5, unit draw seed 1 (\texttt{pipeline/60\_imperfect\_sample.R}), noise seeds
$20260917 + 1000\cdot\mathrm{unit} + \mathrm{replicate}$
(\texttt{pipeline/61\_imperfect\_sweep.R}); Extended Data Fig.~1 and Supplementary Fig.~2, jitter
seed 1; Supplementary Fig.~3A, seed 3 (\texttt{pipeline/51\_harm\_explore\_sample.R});
Fig.~2 bootstrap, seed 1. The main sweep, the knockout integrations,
the boundary scans, the damage-matched search and the treatment arms are deterministic.
The kinetic parameter grids were drawn without a seed (Section~\ref{sec:sampling}).

%% ===========================================================================
\section{Limitations imposed by the data}
\label{sec:limitations}

The following restrictions follow from how the data were generated. Each is stated
where it applies; they are collected here.

\begin{enumerate}
\item \emph{Sterile boundary scans are a subsample.} The $t=0$ boundary scans of
  Supplementary Fig.~1 cover 16{,}990 of the 89{,}328 sterile combinations that passed criterion~(3)
  (19.0\%); the remaining 72{,}338 were never measured. Supplementary Fig.~1B and every sterile number of
  Section~\ref{F-sec:fig5} of the figure methods rest on this random subsample, whose composition matches the
  pool (ON\_mono share, rescue-gap distribution). The pathogenic help pool was measured
  in full.
\item \emph{Fig.~2 uses final states, not attractors.} The knockout experiment computes
  no fixed points, checks no stability and applies no convergence test; a low-burden
  basin is defined by the spread ($>20$ burden units) of the final states at
  $t = 12{,}000$ along the initial-condition line, and a line still moving at that time is
  classified by where it is. In the pathogenic branch the lower-burden final state
  carries more pathogen than the higher-burden one in about half of the systems
  (lower $P_{LP}$ in 49.4\%), so ``lower burden'' there does not mean pathogen
  clearance. Figs.~2F, 2G pool all seven $K_g$ arms of each qualifying unit, and the
  plotted subsampled median hides a wide spread of its per-$K_g$ values
  (Table~\ref{F-tab:fig01-FG} of the figure methods). Each point rests on 13 units per stratum (78 in F, 117
  in G), the size of the smallest ($\bros$, $r_P$) cell after the additional draws.
\item \emph{ON\_mono is a statement about the initial-condition line.} A second stable
  Tregs-ON fixed point exists off the line in 763 of 11{,}159 sterile, 581 of 1{,}270
  pathogenic help and 554 of 1{,}180 pathogenic harm ON\_mono combinations, and three
  pathogenic ON\_mono Tregs-ON trajectories of Extended Data Fig.~3 reach a
  high-burden state. ON\_mono means that no probed initial condition of the line
  reaches a high-burden attractor under Tregs, not that none exists.
\item \emph{The rescue gap has a resolution of one bisection step} (0.078125). A gap of 0
  means a threshold shift smaller than this step; 53.1\% (sterile), 58.7\% (pathogenic
  help) and 65.7\% (pathogenic harm) of the measured gaps are exactly 0.
\item \emph{The Fig.~5 rescue criterion ignores pathogen load.} A run is rescued when
  its tail-mean dead-cell load is below $K_g$; in the pathogenic help branch 258{,}692 of
  367{,}762 rescued runs keep a tail-mean $P_{LP} > 1$.
\item \emph{Missing runs in the Fig.~5 sweep.} 50 complete (unit, replicate) items, all
  pathogenic help (6{,}050 Tregs-ON runs), are absent from the merged output; the stored
  output cannot distinguish an item that never completed from one whose runs were all
  removed by the run filter. Eleven units lose their Tregs-OFF run with them and are
  dropped; the remaining cells carry 7--10 replicates, and every cell of the 440 units
  of Fig.~5 carries 10. The noise jump times are not passed
  to the solver as critical times (checked on a sample with $h_{\max} = 0.1$).
\item \emph{Gate metric.} The gate separates attractors by the total burden; the
  Chebyshev separation would pass 384{,}675 pathogenic combinations instead of 115{,}153.
  The decreasing Tregs-OFF configuration (53{,}881 pathogenic combinations) is excluded.
\item \emph{Pathogenic boundary-scan disagreement.} In 5{,}007 pathogenic combinations
  the Tregs-OFF attractor separation of the boundary stage and of the sweep differ by
  more than 1\% on $P_{LP}$ only, at absolute differences of order $10^{-10}$; they are
  kept in Supplementary Fig.~1C (excluding them changes $\rho$ by at most 0.06).
\item \emph{Fig.~4 proportions} are proportions of the 52{,}010 units run (a uniform
  subsample of a uniform draw from the gated pathogenic combinations in which the pathogen
  persists without Tregs), not frequencies in a population; the units of the draw that the
  cluster tasks did not reach before their time limit are in random order, so their absence
  does not bias the sample.
\item \emph{Treg treatments.} 26 of the 1{,}771 integrated harm combinations are not
  harmed when integrated untreated (their branch was assigned on the fixed points); they
  have no threshold. The numerical robustness of the extreme point of Supplementary Fig.~3B was not re-examined on the final data.
\item \emph{Parameter grids.} The kinetic grids were drawn without a random seed and
  cannot be regenerated; the stored grids are the record. Only grids 1--3 were swept;
  the knockout experiment of Fig.~2 draws from all five grids.
\end{enumerate}

%% ===========================================================================
\section{Software and reproducibility}
\label{sec:software}

\emph{Software.} All computations were done in R. Local runs and all figure renders used
R 4.6.1 (aarch64 macOS) with \texttt{deSolve} 1.42 (integration), \texttt{data.table}
1.18.4, \texttt{ggplot2} 4.0.3, \texttt{patchwork} 1.3.2, \texttt{scales} 1.4.0,
\texttt{ggrepel} 0.9.8, \texttt{MASS} 7.3.65 (\texttt{kde2d}), \texttt{gtable} 0.3.6,
\texttt{viridisLite} 0.4.3, \texttt{ggnewscale} 0.5.2, \texttt{ggh4x} 0.3.1 and the base
packages \texttt{parallel}, \texttt{grid} and \texttt{stats}.
%% from: Rscript -e 'packageVersion(...)' on 2026-10-01; R.version.string; Sys.info()
Cluster stages ran under Slurm on the \texttt{della} cluster, with R taken from a
conda environment (module \texttt{anaconda3/2023.3}); the R and \texttt{deSolve}
versions of that environment were not recorded with the outputs. Re-integrations of
cluster outputs on the local installation agree with the stored values to the
precision given below.

\emph{Computation.} Main sweep: Slurm array of 400 single-core tasks per grid (24 h wall
time, 4 GB per core; \texttt{pipeline/cluster/sweep\_array.sh}).
Efferocytosis knockout: 1{,}764{,}000 integrations, about 19.5 h on 10 cores (about 195
core-hours, or about 67 s per unit per core). Damage-matched experiment (Fig.~4): Slurm array of 400 single-core
tasks (8 h wall time; about 340 integrations and 3.7 core-minutes per unit; every task
reached the wall time), \texttt{pipeline/cluster/damage\_matched\_gated.sh}. Boundary scans of Supplementary Fig.~1: Slurm arrays of 400 single-core tasks
(12 h wall time; about 15 core-minutes per combination); the sterile array exceeded its
wall time, and only its first round was recovered (Section~\ref{F-sec:fig5} of the figure methods).
Stability--specificity sweep: four waves of 400
single-core Slurm tasks (23 h 59 min wall time), 26{,}982 core-hours in total. Treg
treatments: one 32-core node, 130.5 min for the arms and 85.9 min for the thresholds.
Load-balance trajectories: about 2.5 min per branch on 6 cores.

\emph{Stored versus re-run stages.} Expensive integration stages were run once and their
outputs are the data of record; cheap stages were re-run from them, and each expensive
stage was spot-checked by re-integrating a subset on the local installation:
\begin{itemize}
\item \emph{Main sweep} (\texttt{pipeline/01\_sweep.R}, \texttt{02\_assemble.R}): stored.
  The selection stage (\texttt{03\_select.R}) re-run on the stored batches reproduces the
  gated set and the six branch files exactly (\texttt{identical()}). The number of
  assembly merges that joined a pathogen-free fixed point to a pathogen-carrying one was
  not recorded and cannot be reported.
\item \emph{Efferocytosis knockout} (\texttt{06\_effero\_knockout.R}): stored; the unit
  draw re-run (all 10{,}500 units reproduced); three units (504 integrations) re-integrated,
  every final state reproduced to within $3\times10^{-12}$ and the one failed run failing
  again.
\item \emph{Load balance} (\texttt{20\_load\_balance.R}): re-run in full (3{,}360
  trajectories); same combinations and initial conditions, every series of the same
  length, largest absolute difference $2.6\times10^{-4}$ in a stored value.
\item \emph{Boundary scans} (\texttt{21\_corridor\_tzero.R}): stored; two pathogenic
  combinations re-integrated, all 16 (combination, coordinate) rows matching (categorical
  columns, counts and rescue gaps identical; boundaries within $4\times10^{-7}$ absolute
  except on a coordinate whose basin span is $6\times10^{-11}$). Attractor coordinates
  and positions (\texttt{22}, \texttt{23}) re-run, identical. The panel-A boundary sweep
  of Supplementary Fig.~1 is read from its stored cache.
\item \emph{Damage-matched experiment} (\texttt{40\_damage\_matched.R}): stored; 11 units
  covering every search outcome re-run, every numeric column and status identical. The
  stored trajectory subset (\texttt{41}) is the seed-1 draw of 1{,}000 pairs from the current
  filter, so panels B and C use the whole draw. The first peak
  of each arm (\texttt{42}) is stored for all 11{,}940 pairs the filter selects before the
  exclusion of inverted Tregs-ON thresholds, so it covers all 8{,}789 pairs of the current
  filter, and both arms returned a complete integration. The window damage of criterion 3 (\texttt{44}) is stored for
  all 48{,}842 robust pairs.
\item \emph{Stability--specificity sweep} (\texttt{61\_imperfect\_sweep.R}): stored; unit draw
  (\texttt{60}) and rescue probabilities (\texttt{63}) re-run, identical; one (unit,
  replicate) item partly re-integrated (12 of its runs), tail means agreeing to the seventh significant
  digit (largest relative difference $1.5\times10^{-5}$, in a tail variance) and every
  rescue outcome identical. The merge stage (\texttt{62}) needs the per-task output, which
  is held only on the cluster, and was not re-run.
\item \emph{Treg treatments} (\texttt{50\_harm\_rescue\_compute.R}): stored; four
  combinations re-integrated (all 168 arm rows and 12 thresholds identical). The
  help-versus-harm sample of Supplementary Fig.~3A (\texttt{51}) re-run,
  bit-identical.
\end{itemize}
Every figure is rendered from the stored data by its script; apart from the
re-integrations named in the figure sections, no figure script integrates the model.

%% ===========================================================================
\end{document}

